\section{Hyperbolic Geometry and Tilings}

The crystal of experience does not float in featureless space. It lives in a space with a definite shape,
and that shape is the first thing the theory borrows from mathematics. Most of physics quietly assumes flat
space and adds curvature only when gravity forces the issue. The present model does the reverse. It begins
from curvature, because a structure that must grow forever needs room that flat space cannot provide.

This section builds the geometric intuition. We describe the three shapes a space can have, show how a
resident can tell them apart from the inside, settle into hyperbolic space and its Poincare-disk portrait,
explain why such a space keeps opening up while still reading as flat to anyone living in it, and tie the
angle argument back to the $\{5,3,4\}$ crystal at the center of the theory.

\subsection{The three shapes a space can have}

There are exactly three uniform geometries, and all three are already familiar from ordinary surfaces. The
mental move is to stop seeing them as bent sheets inside our world and to see each as a whole world in its
own right.

Take a flat tabletop. Draw a triangle. Its three angles sum to exactly $180$ degrees. Parallel lines hold
the same distance forever. A circle of radius two has exactly twice the rim of a circle of radius one. This
is flat space, the geometry of Euclid. Its curvature is zero.

Now take the surface of a ball. Draw a large triangle from the north pole down to the equator, a quarter of
the way along the equator, then back up to the pole. Every corner is a right angle, so the angles sum to
$270$ degrees, more than $180$. Lines that start out parallel, such as two lines of longitude at the
equator, bend toward each other and meet at the pole. Room runs out. This is round space, positively
curved. A sphere is finite. You can walk all the way around and return home.

Now imagine the opposite. A surface shaped like a saddle, or the ruffled edge of a leaf. Draw a triangle on
it and the angles sum to less than $180$ degrees. Lines that start parallel curve away from each other and
never meet. There is always more room than expected. This is hyperbolic space, negatively curved. It is the
setting for everything that follows.

\begin{table}[h]
\centering
\resizebox{\textwidth}{!}{%
\begin{tabular}{llllll}
\toprule
Geometry & Curvature & Triangle angle sum & Parallels through a point & Ball growth & Regular tilings of the plane \\
\midrule
Spherical & positive & greater than $180$ degrees & none & bounded & finitely many (the Platonic solids) \\
Euclidean & zero & exactly $180$ degrees & exactly one & polynomial & exactly three \\
Hyperbolic & negative & less than $180$ degrees & infinitely many & exponential & infinitely many \\
\bottomrule
\end{tabular}}
\caption{The three uniform geometries and how to tell them apart from inside.}
\label{tab:three-geometries}
\end{table}

\subsection{Telling them apart from inside}

The key point is that a resident need not leave a space to know its curvature. Curvature is measurable from
within, by checking how much room there is. This matters, because we are stuck inside our own universe with
no outside vantage point. Every test below is one a resident can run.

The cleanest test is the angle sum of a triangle. Flat space gives exactly $180$ degrees, round space gives
more, hyperbolic space gives less. The gap from $180$ is proportional to the area of the triangle, so the
larger the triangle, the louder curvature speaks.

The next test is the fate of parallel lines. In flat space they hold their distance forever. In round space
they converge and cross. In hyperbolic space they spread apart and race away from each other.

The deepest test, and the one that drives the theory, is the growth of balls. Stand at a point and ask how
many places sit within one step, within two, within ten. In flat space the room grows steadily, with the
area inside a circle growing as the square of the radius. Double the radius and the room is four times
larger. This is polynomial growth. In round space the room grows and then shrinks back as the space wraps
around on itself. It runs out.

In hyperbolic space the room grows faster and faster, exploding outward. The amount within a given distance
does not merely grow. It grows exponentially. That last fact is the whole reason hyperbolic space matters
here. It has, in a precise sense, endless room near every point. The further you look, the more the world
opens up. The general theory of these growth rates and the geometry of negative curvature is developed by
Cannon and collaborators \cite{cannon1997}.

\subsection{Hyperbolic space and the Poincare disk}

Our eyes were built for flat space, so hyperbolic space feels alien at first. One picture tames it. You
cannot draw the whole of hyperbolic space honestly on a flat page, because there is too much of it, but
there is a faithful way to squeeze all of it into a finite circle. It is the Poincare disk, the device
behind the Escher prints in which figures shrink as they march toward the rim.

Read it as follows. The entire infinite hyperbolic plane is shown inside one round disk. The catch is the
scale. Near the center, distances look normal. Toward the rim, everything is drawn smaller and smaller. A
shape near the edge looks tiny on the page, but in the real hyperbolic world it is exactly the same size as
a shape in the middle. The shrinking is an artifact of cramming infinity into a circle, not a fact about the
space itself.

Two consequences follow. First, the rim is infinitely far away. As you walk toward the edge, your steps are
drawn shorter and shorter on the page, so you never arrive. The edge is not a wall. It is the horizon at
infinity. Second, there is always more room out there. Because room opens up exponentially, the closer you
get to the rim, the more cells fit into each ring. The crowd at the edge outnumbers everything inside it.

In the Poincare disk the shortest path between two points is drawn as an arc that bows gently toward the
center. This looks wrong to a flat-space eye, but it is correct. That bowed arc is the cheapest route. The
straight, cheap paths run through the deep interior, not along the crowded rim. The disk also seems to have
a special center, but that is a lie of the drawing. Hyperbolic space looks exactly the same from every
point and along every direction. You can slide the picture so any point becomes the center, and nothing
about the geometry changes. Our universe has no preferred direction either, and a crystal built in this
directionless space inherits that property for free.

\subsection{The tiny diameter}

Because the count of cells explodes exponentially with distance, the reverse is also true. To reach a vast
number of cells you need travel only a tiny distance. Pack a billion cells around a point and the farthest
of them sits a few steps away, not a billion steps. In flat space a billion cells would force a walk of
roughly a thousand steps out. In hyperbolic space the same billion crowd into a ball of small radius. A
hyperbolic structure of astronomical size can have a diameter that barely grows at all. If everything is
close in graph distance, influence does not have to crawl across a wide expanse. This is the geometric root
of the non-local field that underlies the flat layer we live on.

If space is genuinely saddle-shaped, why does the room around you look flat? For the same reason the Earth
looks flat from a sidewalk. Any curved space looks flat in a small enough patch. Curvature reveals itself
only over large distances or in fine structure. We live in a patch, so day to day space reads as ordinary
and Euclidean.

There is a sharper version of this for hyperbolic space. A horosphere, the limit of a sphere whose center
has drifted off to infinity, has an internal geometry that is perfectly flat. It is a genuine Euclidean
sheet sitting inside a curved world, and the layer we experience as ordinary flat space is exactly this kind
of surface. Curved underneath, flat to the touch.

\subsection{The shape decides the curvature}

A tiling is a way of covering a space with copies of shapes that fit together with no gaps and no overlaps.
The shapes are the tiles, or cells. A tiling is regular when all tiles are the same regular shape and every
vertex looks identical to every other. The same number of tiles meet at each corner, arranged the same way.
A regular tiling has no preferred location and no tunable knobs. Once the cell and the vertex arrangement
are named, the whole infinite pattern is determined.

Stand at a single vertex where tiles meet and add up the corner angles of the tiles that surround it. That
sum is the vertex angle total, and everything follows from comparing it to a full turn of $360$ degrees. If
the angles total exactly $360$ degrees, the tiles lie perfectly flat, and the space is Euclidean. If they
total less than $360$ degrees, there is a wedge of missing angle at each vertex, and to close that gap the
surface must pucker inward into a ball, which is positive curvature. If they total more than $360$ degrees,
there is too much angle to lie flat, and the surface must ruffle outward into a saddle, which is negative
curvature.

The missing or surplus wedge is the angle defect. This is the central move. The angle total at each vertex
is a curvature dial, and because the angle total is fixed entirely by the tile shape and how many tiles
meet, the tile choice and the curvature are one decision, not two.

The flat plane allows exactly three regular tilings.
\begin{itemize}
\item Triangles, six to a corner.
\item Squares, four to a corner.
\item Hexagons, three to a corner.
\end{itemize}
That is the complete list. The reason is arithmetic. Squares contribute $90$ degrees each, so four make
$360$. Triangles contribute $60$, so six make $360$. Hexagons contribute $120$, so three make $360$.
Pentagons contribute $108$, and no whole number of these lands on $360$, so the pentagon is exiled from the
flat plane.

Two everyday examples sharpen the picture. The checkerboard is the square tiling with its tiles two-colored
in alternation. The coloring changes nothing about the geometry, but it reveals a bipartite structure in
which every tile borders only tiles of the opposite color. The soccerball, the truncated icosahedron of
twelve pentagons and twenty hexagons fitted onto a sphere, is the shape of the $\mathrm{C}_{60}$ molecule,
and it carries the same fivefold symmetry and golden ratio that the theory's dodecahedral cell does.

\subsection{Why hyperbolic space is abundant}

The flat plane offers three regular tilings. The sphere offers a small handful, the Platonic solids seen as
tilings of a round surface. Hyperbolic space offers infinitely many. The reason is the surplus. Hyperbolic
space has room to spare, so any arrangement whose angle total comes out over $360$ is welcome, because the
surplus is exactly what the saddle geometry absorbs. Want seven-sided tiles, three to a corner? Their angles
overflow $360$, so flat space refuses and hyperbolic space accepts. Want hundred-sided tiles? Also fine. The
overflow that flat and round space cannot tolerate is the very thing the saddle needs.

This is the property the theory depends on, because we want two things at once. We want perfect regularity,
every cell and every vertex identical with no special points. We want endless room, a crystal that grows
forever without folding back into a small finite shape. The sphere gives regularity but no room. The flat
plane gives room but almost no regular patterns. Only hyperbolic space gives regularity and endless room
together.

So the crystal of experience must be a regular tiling of hyperbolic space, and the abundance of such tilings
is what makes the precise, exotic cell the theory needs available at all. The same angle-defect logic
carries to three dimensions, where cells pack around an edge and the relevant angle is the cell's dihedral
angle. That is the route to the $\{5,3,4\}$ honeycomb, taken up in the next section.
